\section{Evidence required for institutional adoption}
\label{exec17:signed-adoption}

The evaluation object must identify what was assessed. A proposed release
manifest binds the reader PDF and source audit; canonical schemas and operation
registry; runtime and adapter builds; dependency and protocol versions; the
consumer decoder; the scientific target and sampling law; the declared
platform; and the acceptance criteria. A change to any item starts an impact
review. Prior receipts may be carried forward only when their owner records
which claims remain applicable and why. A hash establishes byte identity, not
who signed those bytes, whether that person had authority, or whether the
scientific statement is true.

The exact authored gate table contains six assumed conditions: independent
scientific review; authorized bank data and truth; actual consumer acceptance;
target-platform acceptance; a matched economic case; and a sponsor decision.
There are $2^6=64$ Boolean combinations, but only the all-true combination is
eligible for scoped adoption under the stated conjunction. This is a logic
check on a proposed rule. It supplies zero authentic receipts, identifies no
signer, and performs no cryptographic verification. Institutional acceptance additionally requires identifiable responsible
owners and authentic evidence for scientific review, bank truth, platform,
consumers, workflow economics and the sponsor decision.~\citep{exec17_continuity_reference,exec17_continuity_contract}

Local preparation can still be complete enough to hand over. The package
should contain an exact-release manifest, a criterion-by-criterion evaluation
protocol, a custody-safe bank data manifest, the consumer and platform test
plans, a full-cost ledger, a conflicts statement, and a receipt schema with
review scope, owner role and organization, evidence identity, criteria,
unresolved findings, and assessed release. The intake script can verify
required fields and hashes. It explicitly cannot authenticate a signature,
establish legal access, validate bank truth, or decide scientific adequacy.
Those determinations belong to the authentic responsible owner and reviewer.
The existing external-input contract states this boundary directly.
\citep{exec17_continuity_contract}

The six-field release example changes the reader, core, adapter, consumer,
scientific context, or platform one at a time. Every change breaks exact
release equality and triggers impact review; it does not automatically erase
all prior evidence. A scoped owner may establish that an earlier test remains
applicable, but the evaluator must bind that judgment to the successor
release. Otherwise the stale receipt is not evidence for the changed system.
The proposal deliberately has no success branch that manufactures the missing
owner, signature, bank history, or institutional decision.
\citep{exec17_continuity_reference,exec17_continuity_contract}

A prototype comparison can test whether the same scientific context survives
library calls, command-line decoding, protocol transport and declared job or
storage failures. Its record needs full request, response, cost and refusal
denominators. Success under constructed failures supports a scoped consumer
and platform trial; it supplies no authentic bank truth or institutional
approval. The economic decision to fund that trial depends on its information
value and a meaningful precision requirement, as developed in
Chapter~\ref{rewrite:investment-decision}.
\label{exec17:integration-continuity-end}
